\documentclass{report}

\usepackage{fontspec}
\usepackage{adjustbox}
\usepackage{rotating}
\usepackage{polyglossia}
\usepackage{float}
\usepackage{amsmath}
\usepackage{amssymb}
\usepackage{makebox}
\usepackage{multirow}


\setdefaultlanguage{ukrainian}
\setotherlanguages{english}

\setmainfont{CMU Serif}
\newfontfamily\cyrillicfont{CMU Serif}

\usepackage{listings}

\usepackage{geometry}
\geometry{
    a4paper,
    left = 30mm,
    top = 20mm,
    right = 10mm,
    bottom = 20mm
}

\lstset{
    basicstyle=\ttfamily\small,
    frame=single,
    breaklines=true,
    numbers=left,
    numberstyle=\tiny,
    captionpos=b,
}

\begin{document}

\pagestyle{empty}

\begin{titlepage}
    \begin{center}

        \textbf{НАЦІОНАЛЬНИЙ ТЕХНІЧНИЙ УНІВЕРСИТЕТ УКРАЇНИ}\\
        “КИЇВСЬКИЙ ПОЛІТЕХНІЧНИЙ ІНСТИТУТ ім. Ігоря Сікорського”\\
        Навчально-науковий фізико-технічний інститут\\
        Кафедра математичних методів захисту інформації

        \vspace{6cm}

        \Large \textbf{Методи криптоаналізу. Частина 1}\\
        Комп'ютерний практикум №2\\

    \end{center}

    \vspace{10cm}
    \begin{flushright}
        Виконали:\\
        Студенти групи ФІ-52мн\\
        Ісаченко Нікіта\\
        Фурутіна Євгенія
    \end{flushright}

    \vspace*{2.5cm}

    \begin{center}
        2026 р.
    \end{center}
\end{titlepage}

\begin{table}[h!]
\centering
\caption{Спотворення за допомогою шифру Віженера ($r = 5$)}
\label{tab:vigenere_r5}
\renewcommand{\arraystretch}{1.3}
\begin{tabular}{|c|c||c|c||c|c|}
\hline
\multicolumn{2}{|c||}{} &
\multicolumn{2}{c||}{$l = 1$} &
\multicolumn{2}{c|}{$l = 2$} \\
\cline{3-6}
\multicolumn{2}{|c||}{} &
FP & FN & FP & FN \\
\hline
\multirow{5}{*}{10}
 & 2.0 (обрані порогові значення) &  &  &  &  \\
 & 2.1 &  &  &  &  \\
 & 2.2 &  &  &  &  \\
 & 2.3 &  &  &  &  \\
 & 4.0 &  &  &  &  \\
\hline
\end{tabular}
\end{table}

\end{document}